% ---------------------------------------------------------------------------
% Pillar 2: Zero-Variance Fraction diagnostic -- Iter 90
% Recovery-Asymmetry: classify each ZVF-alarm episode as TRANSIENT
% (recovered within K=5 steps) vs SUSTAINED, and ask whether
% post-recovery heldout-acc matches post-sustained heldout-acc.
% Materialised from scripts/zvf_iter90.py:
%   experiments/results/zvf_iter90_episodes.tsv       (54 rows)
%   experiments/results/zvf_iter90_post_episode.tsv   (27 rows)
%   experiments/results/zvf_iter90_recovery.tsv       (9 rows)
%   experiments/results/zvf_iter90_summary.tsv       (15 KV pairs)
%   figures/zvf_iter90_recovery.{pdf,png}
% ---------------------------------------------------------------------------

\subsection{ZVF Recovery Asymmetry: Transient Alarms Do Not Rescue Accuracy}
\label{sec:zvf-recovery}

Iter~86 closed the decision-theoretic loop on the ZVF alarm by
selecting $(\tau=0.6,\, K=5)$ as the optimal cost-vs-savings pair on
the variance-mitigation pool (\texttt{zvf\_iter86\_optimal\_tau.tsv}).
That framing implicitly assumed the alarm is a sharp signal:
\emph{stop-on-alarm} is conservative because alarms lead to collapse
within $\le 7$ steps with certainty. But the framing leaves a clean
question for any practitioner who hesitates to stop training:
\emph{how likely is the alarm itself to recover before the collapse
arrives?}

Iter~90 answers this. For every variance-mitigation trace, we
detect contiguous alarm episodes
$\mathcal{E}_i = \{t : \mathrm{ZVF}_{i,t} > \tau\}$ with
$\tau=0.5$ fixed to the iter82 hazard threshold. Each episode is
labelled \emph{transient} if ZVF drops back to $\le 0.5$ within
$K=5$ steps of the alarm's last above-threshold sample (matching
iter86's $k_{\text{persist}}$), and \emph{sustained} otherwise.

\paragraph{Recovery rates are bimodal across methods.}
On 45 traces (9 methods $\times$ 5 seeds), we observe 54 alarm
episodes, of which 21 transient (pooled recovery rate
$0.389$) and 33 sustained ($0.611$).
Method-stratified rates
(\texttt{zvf\_iter90\_recovery.tsv}) split sharply:
%
\begin{itemize}\itemsep1pt
  \item \emph{Zero-recovery methods} (rate=0):
        \textsc{grpo} (5/5 sustained).
        \textsc{scafgrpo} and \textsc{es} never alarm at all
        (their ZVF stays $\le 0.5$ across the 123-step trajectory),
        so the recovery-rate denominator is vacuously zero.
  \item \emph{Low-recovery methods} (rate $\in [0.17, 0.38]$):
        \textsc{aero} 0.167,
        \textsc{cppo} 0.286, \textsc{ngrpo} 0.286,
        \textsc{mcgrpo} 0.375.
  \item \emph{High-recovery methods} (rate $\ge 0.60$):
        \textsc{areal} 0.600, \textsc{gift} 0.636.
\end{itemize}
%
Vanilla \textsc{grpo}, which iter82 already flagged as having the
highest in-alarm hazard ratio ($\mathrm{HR}=143$), now sharpens to
\emph{zero recovery}: any ZVF alarm in pure GRPO is permanent within
the variance-mitigation budget. Conversely \textsc{gift} recovers in
two of three alarms.

\paragraph{Transient alarms do \emph{not} rescue accuracy.}
The sharper finding is what recovery does \emph{not} buy. Comparing
the post-episode $15$-step trailing heldout-acc
(\texttt{zvf\_iter90\_post\_episode.tsv}):
%
\begin{equation*}
  \begin{array}{lccc}
    \text{method} & \overline{ha}_{\text{recovered}} &
    \overline{ha}_{\text{sustained}} & \Delta \\[2pt]
    \textsc{aero}  & 0.389 & 0.405 & -0.017 \\
    \textsc{cppo}  & 0.367 & 0.396 & -0.029 \\
    \textsc{ngrpo} & 0.360 & 0.393 & -0.033 \\
    \textsc{mcgrpo}& 0.395 & 0.403 & -0.007 \\
    \textsc{gift}  & 0.405 & 0.409 & -0.004 \\
    \textsc{areal} & 0.406 & 0.406 & \phantom{-}0.000 \\
  \end{array}
\end{equation*}
%
Across every method, the post-recovery heldout-acc is at most
$\Delta=0.033$ \emph{lower} than the post-sustained heldout-acc, and
two methods (\textsc{areal}, \textsc{gift}) have $\Delta < 0.01$.
\emph{ZVF recovery and accuracy recovery are decoupled.} Even though
the ZVF signal returns to the safe band, the heldout accuracy does
not rebound; it converges to the same plateau as the sustained
trajectory.
%
\paragraph{False-friend correction to iter86.}
Iter~86 recommends \emph{stop-on-alarm} at $\tau=0.6$ on the
implicit assumption that \emph{silent ZVF = healthy gradient}.
Iter~90 falsifies the dual assumption: \emph{recovered ZVF
= restored gradient}. A trainer that waits $K=5$ post-alarm steps
to see if ZVF drops back, hoping to resume training only when the
alarm is transient, would resume into the same accuracy plateau that
an immediate stop would have avoided. The only methods where
transient-vs-sustained is meaningfully distinguishable on accuracy
are \textsc{cppo} and \textsc{ngrpo}, and even there the gap is
only $\approx 3$ percentage points.

\paragraph{Method-level prediction.}
GRPO's $0\%$ recovery combined with iter82's $\mathrm{HR}=143$
gives the cleanest operational rule: any GRPO alarm within the
variance-mitigation budget is both necessary ($p_{\text{cov}}=1.0$)
\emph{and} sufficient ($p_{\text{recov|alarm}}=0.0$) for collapse,
so \emph{GRPO alarms are decision-complete}. Methods in the
high-recovery tier (\textsc{gift}, \textsc{areal}) have lower
in-alarm hazard than GRPO even when the alarm fires, so stopping
on every alarm there is the costlier half of iter86's pairing.

\paragraph{Deliverables.} \texttt{scripts/zvf\_iter90.py} (~190 LOC,
stdlib $\cup$ matplotlib); \texttt{zvf\_iter90\_episodes.tsv}
(54 alarm episodes with start/end/length/recovery label);
\texttt{zvf\_iter90\_post\_episode.tsv} (27 (method, post-cat) rows);
\texttt{zvf\_iter90\_recovery.tsv} (9 methods);
\texttt{zvf\_iter90\_summary.tsv} / \texttt{.json};
\texttt{figures/zvf\_iter90\_recovery.pdf} ($\sim 5$~KB).
